\documentclass{amsart}
\usepackage{amsmath,amssymb,amsthm,hyperref}
\newtheorem{theorem}{Theorem}
\newtheorem{lemma}{Lemma}
\newtheorem{proposition}{Proposition}
\newtheorem{corollary}{Corollary}
\theoremstyle{definition}
\newtheorem{remark}{Remark}
\newcommand{\Gr}{\mathrm{Gr}}
\newcommand{\Hilb}{\mathrm{Hilb}}
\newcommand{\qbin}[2]{\genfrac{[}{]}{0pt}{}{#1}{#2}_q}
\newcommand{\Q}{\mathbb{Q}}
\newcommand{\rkl}{R^{k,\ell}}
\newcommand{\Rklm}{R^{k,\ell,m}}
\begin{document}

\title{The Hilbert series of the successive subalgebras
$\Q\langle e_1,\dots ,e_m\rangle$ of Grassmannian cohomology}
\maketitle

Throughout, $\Lambda=\Lambda_{\Q}$ is the ring of symmetric functions over $\Q$,
$s_\lambda$ the Schur functions, $e_i$ the elementary and $h_i$ the complete
homogeneous symmetric functions, all graded by weight $|\lambda|=\sum_i\lambda_i$.
For a partition $\lambda$ we write $\ell(\lambda)$ for its number of (positive)
parts, $\lambda_1$ for its largest part, $\lambda'$ for the conjugate partition
(so $\lambda'_1=\ell(\lambda)$ is the longest column height), and
$\lambda\subseteq k\times\ell$ when
$\ell(\lambda)\le k$ and $\lambda_1\le\ell$. We use the standard $q$-integers
$[n]_q=\frac{1-q^n}{1-q}$, $[n]_q!=\prod_{i=1}^n[i]_q$, and Gaussian binomial
coefficients $\qbin{n}{k}=\frac{[n]_q!}{[k]_q![n-k]_q!}$.

\medskip
\noindent\textbf{The ring.} As recalled in the statement,
\[
\rkl=H^{*}(\Gr(k,\mathbb{C}^{k+\ell});\Q)\cong
\Lambda\big/ I^{k,\ell},
\qquad
I^{k,\ell}=\big\langle\, s_\lambda \ :\ \lambda\not\subseteq k\times\ell\,\big\rangle ,
\]
and a $\Q$-basis of $\rkl$ is $\{\,\overline{s_\lambda}\ :\ \lambda\subseteq k\times\ell\,\}$;
hence
\begin{equation}\label{eq:HilbR}
\Hilb(\rkl;q)=\sum_{\lambda\subseteq k\times\ell}q^{|\lambda|}=\qbin{k+\ell}{k}.
\end{equation}
For $0\le m\le\ell$ we set $\Rklm=\Q\langle e_1,\dots,e_m\rangle\subseteq\rkl$
(with $R^{k,\ell,0}=\Q$), a graded subalgebra, and
\[
Q_{k,\ell,m}(q):=\Hilb\!\left(\Rklm/R^{k,\ell,m-1};q\right)
=\Hilb(\Rklm;q)-\Hilb(R^{k,\ell,m-1};q),
\]
where the last equality holds because $R^{k,\ell,m-1}\subseteq\Rklm$ are graded
subspaces, so each graded piece of the quotient has dimension the difference of
dimensions. We must prove, for all $k\ge1$ and $1\le m\le\ell$,
\begin{equation}\label{eq:main}
Q_{k,\ell,m}(q)=q^{m}\,\qbin{\ell}{m}\,
\sum_{j=0}^{k-m} q^{\,j(\ell-m+1)}\,\qbin{m+j-1}{j}
=:F_{k,\ell,m}(q).
\end{equation}

\bigskip
The proof proceeds by induction on $k$ (for fixed $\ell$ and $m$). The structural
input is a surjection $\pi\colon\rkl\twoheadrightarrow R^{k-1,\ell}$ which restricts
to the subalgebras and whose kernel is principal. We first record elementary
reductions and the behaviour of the right-hand side, then develop the structural
lemmas, and finally combine them.

\section{Elementary reductions and the right-hand side}

\begin{lemma}\label{lem:trivial}
If $m\ge k$ then $\Rklm=\rkl$. In particular $R^{k,\ell,\ell}=\rkl$, and the
nontrivial range of \eqref{eq:main} is $1\le m\le\min(k,\ell)$; the cases
$\ell\ge m>k$ do not occur because $m\le\ell$ and we take $k\ge m$ in the
induction (cases $k<m$ being subsumed by $\Rklm=\rkl$).
\end{lemma}

\begin{proof}
The ring $\rkl$ is generated by $e_1,\dots,e_k$. If $m\ge k$ then
$e_1,\dots,e_m$ already contain all generators, so $\Rklm=\rkl$.
For $R^{k,\ell,\ell}=\rkl$ when $k>\ell$: the generating-function identity
$\big(\sum_{i\ge0}(-1)^i e_i z^i\big)\big(\sum_{j\ge0}h_j z^j\big)=1$ gives, for
$n\ge1$,
\[
\sum_{i=0}^{n}(-1)^i e_i\,h_{n-i}=0 .
\]
For $1\le j\le\ell$, induction on $j$ via this relation (with $n=j$) expresses
$h_j$ as a polynomial in $e_1,\dots,e_j$. For $n=\ell+1,\dots,k$, using
$h_{\ell+1}=\dots=h_{\ell+k}=0$ in $\rkl$ and that $h_1,\dots,h_\ell$ are
polynomials in $e_1,\dots,e_\ell$, the relation with this $n$ solves $(-1)^n e_n$
(its top term, since $h_0=1$) in terms of $e_0,\dots,e_{n-1}$ and $h_1,\dots,h_\ell$,
hence in terms of $e_1,\dots,e_{n-1}$. By induction on $n$ each
$e_{\ell+1},\dots,e_k$ lies in $\Q\langle e_1,\dots,e_\ell\rangle$, so
$R^{k,\ell,\ell}=\rkl$.
\end{proof}

\begin{lemma}[Recursion for the right-hand side]\label{lem:RHSrec}
For all $1\le m\le\ell$ and $k>m$,
\begin{equation}\label{eq:RHSrec}
F_{k,\ell,m}(q)-F_{k-1,\ell,m}(q)
= q^{\,m+(k-m)(\ell-m+1)}\,\qbin{\ell}{m}\,\qbin{k-1}{m-1},
\end{equation}
and the base value is
\begin{equation}\label{eq:RHSbase}
F_{m,\ell,m}(q)=q^{m}\,\qbin{\ell}{m}.
\end{equation}
\end{lemma}

\begin{proof}
In $F_{k,\ell,m}$ the summand $\qbin{m+j-1}{j}$ does not depend on $k$, and the
$j$-range is $0\le j\le k-m$. Increasing $k$ by one only appends the single term
$j=k-m$:
\[
F_{k,\ell,m}-F_{k-1,\ell,m}
= q^{m}\qbin{\ell}{m}\,q^{(k-m)(\ell-m+1)}\qbin{m+(k-m)-1}{\,k-m\,}.
\]
Now $\qbin{k-1}{k-m}=\qbin{k-1}{m-1}$ by the symmetry $\qbin{n}{a}=\qbin{n}{n-a}$,
which gives \eqref{eq:RHSrec}. For \eqref{eq:RHSbase}, when $k=m$ the $j$-range is
the single value $j=0$, with $q^{0}\qbin{m-1}{0}=1$, so
$F_{m,\ell,m}=q^{m}\qbin{\ell}{m}$.
\end{proof}

Thus it suffices to prove that the left-hand side $Q_{k,\ell,m}$ obeys the same
recursion and base value:
\begin{align}
\label{eq:LHSrec}
Q_{k,\ell,m}(q)-Q_{k-1,\ell,m}(q)
&= q^{\,m+(k-m)(\ell-m+1)}\,\qbin{\ell}{m}\,\qbin{k-1}{m-1},
&&(k>m),\\
\label{eq:LHSbase}
Q_{m,\ell,m}(q)&=q^{m}\,\qbin{\ell}{m}.
\end{align}

\begin{remark}[Verification]\label{rem:verif}
Both \eqref{eq:main}, and the equivalent pair \eqref{eq:LHSrec}--\eqref{eq:LHSbase},
were checked symbolically for all $1\le k,\ell\le5$ and $1\le m\le\ell$ by computing
$\Rklm$ inside the Schur basis of $\rkl$ (multiplication by $e_i$ is the Pieri rule
for vertical strips, truncated to $k\times\ell$) and row-reducing over $\Q$; all
cases agree. These computations also confirm the structural objects used below
(the surjection $\pi$, its kernel $e_k\rkl$, and the base-case dimension count).
\end{remark}

\section{The column-height span $F_m$ is \emph{not} the subalgebra}
\label{sec:Fm}

\noindent\emph{Resolution of the column-height basis question.} The lemma proposed
for this part---``$F_m=\Rklm$, so $\Rklm/R^{k,\ell,m-1}$ has Schur-class basis
$\{\overline{s_\lambda}:\lambda\subseteq k\times\ell,\ \lambda'_1=m\}$''---is
\emph{false}. We prove the following negative result, which both refutes that lemma
and shows that \emph{no} basis of $\Rklm$ can be cut out by a column-height (or any
monotone shape) threshold; the correct count $Q_{k,\ell,m}$ is obtained instead by
the surjection $\pi$ of Section~\ref{sec:Fm}'s sequel.

\begin{proposition}[The column-height basis does not exist]\label{prop:nocolbasis}
Let $1\le m\le\ell$. With
$F_m=\operatorname{span}_\Q\{\overline{s_\lambda}:\lambda\subseteq k\times\ell,\
\lambda'_1\le m\}$ as below, the following hold.
\begin{enumerate}
\item[\textup{(a)}] $\Rklm\not\subseteq F_m$ whenever $2m\le k$ \textup{(}e.g.\
$(k,\ell,m)=(2,1,1)$\textup{)}: $\overline{e_m^{\,2}}\in\Rklm$ has a nonzero
$\overline{s_{(1^{2m})}}$-component and $(1^{2m})'_1=2m>m$.
\item[\textup{(b)}] $F_m\not\subseteq\Rklm$ in general \textup{(}e.g.\
$(k,\ell,m)=(2,2,1)$\textup{)}: $\overline{s_{(2)}}\in F_1$ but
$\overline{s_{(2)}}\notin R^{2,2,1}=\Q\langle e_1\rangle$.
\item[\textup{(c)}] Consequently $\{\overline{s_\lambda}:\lambda'_1=m\}$ is
neither contained in nor spans $\Rklm/R^{k,\ell,m-1}$, and
$\sum_{\lambda'_1=m}q^{|\lambda|}\ne Q_{k,\ell,m}(q)$ in general; e.g.\ for
$(k,\ell,m)=(3,3,2)$, $\dim_\Q\Rklm=19$ while $\dim_\Q F_m=10$.
\end{enumerate}
In particular $\Rklm$ has no Schur-class basis selected by a column-height bound.
\end{proposition}

The detailed verifications of \textup{(a)}--\textup{(c)} are
Lemma~\ref{lem:Fmfails}, Remark~\ref{rem:Fmnumerics}, and
Remark~\ref{rem:whatswrong} below; all stated numerics were confirmed by direct
computation in the Schur basis of $\rkl$.

It is tempting to look for a Schur-class basis of $\Rklm$ cut out by a column-height
condition. Concretely, define the column-height span
\[
F_m \;:=\; \operatorname{span}_{\Q}\bigl\{\,\overline{s_\lambda}\ :\
\lambda\subseteq k\times\ell,\ \lambda'_1\le m\,\bigr\}\subseteq\rkl ,
\]
i.e.\ the span of those Schur classes whose Young diagram has no column taller than
$m$. One might hope that $F_m=\Rklm$, which would give the clean basis
$\{\overline{s_\lambda}:\lambda\subseteq k\times\ell,\ \lambda'_1=m\}$ of
$\Rklm/R^{k,\ell,m-1}$ with Hilbert series $\sum_{\lambda'_1=m}q^{|\lambda|}$. We
record that \emph{this is false}, and isolate exactly why, so that no later argument
relies on it.

\begin{lemma}[The column-height span fails]\label{lem:Fmfails}
There exist $k,\ell,m$ with $1\le m\le\ell$ for which $F_m\neq\Rklm$. Indeed:
\begin{enumerate}
\item[(i)] $\Rklm\not\subseteq F_m$ already at $(k,\ell,m)=(2,1,1)$: there
$\overline{e_1^{\,2}}=\overline{s_{(1,1)}}\in R^{2,1,1}$, yet $(1,1)'_1=2>1$, so
$\overline{s_{(1,1)}}\notin F_1$.
\item[(ii)] More generally, for every $m\ge1$ and $\ell\ge m$ one has
$\overline{e_m^{\,2}}\in\Rklm$ with a nonzero $\overline{s_{(1^{2m})}}$-component
whenever $2m\le k$, and $(1^{2m})'_1=2m>m$; hence $\Rklm\not\subseteq F_m$ for all
such $(k,\ell,m)$.
\end{enumerate}
\end{lemma}

\begin{proof}
The point is that the generators $e_i=s_{(1^i)}$ are themselves \emph{tall} classes,
and the vertical-strip Pieri rule \emph{increases} column heights: for a partition
$\mu$ with $\ell(\mu)\le k$,
\[
e_i\cdot s_\mu=\sum_{\mu\,\subset\,\nu}s_\nu ,
\]
the sum over $\nu$ with $\nu/\mu$ a vertical $i$-strip (at most one new box per row),
so $\nu'_1=\ell(\nu)$ can exceed $\mu'_1=\ell(\mu)$. Thus a product of the $e_i$
($i\le m$) need not lie in $F_m$.

(i) In $\Lambda$, $e_1^2=h_1 e_1=s_{(2)}+s_{(1,1)}$ by the Pieri rule (the two
addable single boxes to the single cell $(1)$). In $R^{2,1}$ we have $\ell=1$, so
$\overline{s_{(2)}}=0$ (since $2>\ell$), leaving
$\overline{e_1^{\,2}}=\overline{s_{(1,1)}}\neq0$ (since $(1,1)\subseteq2\times1$).
Hence $\overline{s_{(1,1)}}\in R^{2,1,1}$. Its column height is $(1,1)'_1=2>1$, so
$\overline{s_{(1,1)}}\notin F_1$. Therefore $R^{2,1,1}\not\subseteq F_1$.

(ii) In $\Lambda$, $e_m^2=\sum_\nu s_\nu$ summed over $\nu$ with $\nu/(1^m)$ a
vertical $m$-strip. One such $\nu$ is $(1^{2m})$ (stack a second full column of
height $m$ under the first, which is a vertical $m$-strip), and the corresponding
Littlewood--Richardson coefficient is $1$; moreover $(1^{2m})$ is the unique term of
$e_m^2$ of shape $(1^{2m})$, so its coefficient in $e_m^2$ is exactly $1$. If
$2m\le k$ \emph{(}so that $(1^{2m})\subseteq k\times\ell$, as $\ell\ge m\ge1$
ensures the part-size bound\emph{)} then $\overline{s_{(1^{2m})}}\neq0$ in $\rkl$, and it is a Schur
basis vector distinct from all other shapes appearing in $\overline{e_m^{\,2}}$.
Hence the expansion of $\overline{e_m^{\,2}}$ in the Schur basis of $\rkl$ has a
nonzero $\overline{s_{(1^{2m})}}$-coefficient. Since $(1^{2m})'_1=2m>m$, the class
$\overline{e_m^{\,2}}\in\Rklm$ is not a $\Q$-combination of Schur classes of column
height $\le m$, i.e.\ $\overline{e_m^{\,2}}\notin F_m$. Thus
$\Rklm\not\subseteq F_m$.
\end{proof}

\begin{remark}[Numerical scope of the failure]\label{rem:Fmnumerics}
A direct computation of $\Hilb(F_m;q)=\sum_{\lambda\subseteq k\times\ell,\
\lambda'_1\le m}q^{|\lambda|}$ and of $\Hilb(\Rklm;q)$ (by closing $\{1\}$ under
multiplication by $e_1,\dots,e_m$ in the Schur basis of $\rkl$ and row-reducing over
$\Q$) shows $\Hilb(F_m;q)\neq\Hilb(\Rklm;q)$ for every $(k,\ell,m)$ with
$1\le m<\ell$ and $m<k$ in the range $1\le k,\ell\le4$; moreover $F_m$ and $\Rklm$
are \emph{incomparable} subspaces (neither contains the other), as shown in
Lemma~\ref{lem:Fmfails} and Remark~\ref{rem:whatswrong}. For instance, for
$(k,\ell,m)=(3,3,2)$ one finds $\dim\Rklm=19$ but $\dim F_2=10$, while
$\dim(R^{3,3})=20$. Consequently the would-be quotient basis
$\{\overline{s_\lambda}:\lambda'_1=m\}$ has the wrong cardinality, and its naive
Hilbert series $\sum_{\lambda'_1=m}q^{|\lambda|}$ does \emph{not} equal
$Q_{k,\ell,m}(q)$ in general. We have separately verified that the correct value
$Q_{k,\ell,m}(q)$ \emph{does} agree with the target $F_{k,\ell,m}(q)$ of
\eqref{eq:main} for all $1\le k,\ell\le5$, $1\le m\le\ell$ (Remark~\ref{rem:verif});
it is only the column-height \emph{characterization of a basis} that is incorrect.
\end{remark}

\begin{remark}[Both proposed inclusions fail; no shape-statistic basis exists]\label{rem:whatswrong}
\emph{Neither} inclusion in the proposed argument holds.

\emph{$\Rklm\subseteq F_m$ is false} (Lemma~\ref{lem:Fmfails}): multiplying by
$e_i=s_{(1^i)}$ adds a \emph{vertical} strip, which \emph{raises} column heights, so a
product $e_{i_1}\cdots e_{i_s}$ generically has Schur components of column height up to
$\sum_t i_t$, exceeding $m$. The claim that ``box truncation only deletes terms and so
preserves a column bound'' is irrelevant: the tall terms are present \emph{before}
truncation and need not be truncated, e.g.\ $(1,1)\subseteq2\times1$ in
Lemma~\ref{lem:Fmfails}(i).

\emph{$F_m\subseteq\Rklm$ is also false.} The dual Jacobi--Trudi identity
$s_\lambda=\det\bigl(e_{\lambda'_i-i+j}\bigr)_{1\le i,j\le\lambda_1}$ expresses
$s_\lambda$ as a polynomial in $e_1,\dots,e_{N(\lambda)}$ with
$N(\lambda)=\max_{i,j}(\lambda'_i-i+j)=\lambda'_1+\lambda_1-1=\ell(\lambda)+\lambda_1-1$
(the maximum at $i=1$, $j=\lambda_1$): the required generators are controlled by
\emph{both} the column height and the row width, not by the column height alone.
A clean counterexample is $\lambda=(2)$: its longest column has height $1$, so
$\overline{s_{(2)}}\in F_1$; but $s_{(2)}=e_1^{\,2}-e_2$ involves $e_2$, and indeed
$\overline{s_{(2)}}\notin R^{2,2,1}=\Q\langle e_1\rangle$. Concretely in $R^{2,2}$ the
degree-$2$ part of $R^{2,2,1}$ is the line
$\Q\cdot\overline{e_1^{\,2}}=\Q\cdot(\overline{s_{(2)}}+\overline{s_{(1,1)}})$, which
does not contain $\overline{s_{(2)}}$ (as $\overline{s_{(2)}},\overline{s_{(1,1)}}$ are
distinct basis vectors of $R^{2,2}$). Hence $F_1\not\subseteq R^{2,2,1}$.

\emph{Moreover, no monotone shape statistic detects membership.} A direct computation
of the minimal $m$ with $\overline{s_\lambda}\in\Rklm$ (for $\lambda\subseteq
k\times\ell$, $1\le k,\ell\le4$) shows it is not a function of $\lambda_1$, of
$\ell(\lambda)=\lambda'_1$, or of $|\lambda|$ alone; e.g.\ in $R^{4,4}$ one has
$\overline{s_{(1,1)}}\in\Rklm$ only for $m\ge2$, whereas
$\overline{s_{(4,4,4,3)}}\in\Rklm$ already for $m\ge1$. Consequently $\Rklm$ admits
\emph{no} Schur-class basis cut out by a column-height (or row-width, or size)
threshold, and $Q_{k,\ell,m}$ cannot be read off from any such shape condition. The
proof must instead proceed by the comparison surjection $\pi$ of the next section,
which avoids any explicit Schur or monomial basis of the subalgebra.
\end{remark}

\section{The comparison surjection and its kernel}

\begin{lemma}[Comparison surjection]\label{lem:pi}
For $k\ge2$ there is a surjective graded $\Q$-algebra homomorphism
\[
\pi\colon \rkl\twoheadrightarrow R^{k-1,\ell},\qquad \pi(e_i)=e_i\ (1\le i\le k-1),
\quad \pi(e_k)=0,
\]
induced on quotients by $\mathrm{id}_\Lambda$, using $I^{k,\ell}\subseteq I^{k-1,\ell}$.
Its kernel is the principal ideal
\begin{equation}\label{eq:kerpi}
\ker\pi=e_k\cdot\rkl
=\operatorname{span}_{\Q}\{\,\overline{s_\lambda}\ :\ \lambda\subseteq k\times\ell,\ \ell(\lambda)=k\,\},
\end{equation}
and as a graded vector space
\begin{equation}\label{eq:kerHilb}
\Hilb(\ker\pi;q)=q^{k}\,\qbin{k-1+\ell}{k}.
\end{equation}
\end{lemma}

\begin{proof}
Since $(k-1)\times\ell\subseteq k\times\ell$, every $\lambda\not\subseteq k\times\ell$
also satisfies $\lambda\not\subseteq(k-1)\times\ell$, so $I^{k,\ell}\subseteq
I^{k-1,\ell}$ as ideals of $\Lambda$. Hence the identity of $\Lambda$ descends to a
surjective graded algebra map $\pi\colon\Lambda/I^{k,\ell}\to\Lambda/I^{k-1,\ell}$,
i.e.\ $\pi\colon\rkl\twoheadrightarrow R^{k-1,\ell}$. As $e_i$ maps to $e_i$ in
$\Lambda$, we get $\pi(\overline{e_i})=\overline{e_i}$; and
$\overline{e_k}=\overline{s_{(1^k)}}=0$ in $R^{k-1,\ell}$ because $(1^k)$ has $k>k-1$
rows, i.e.\ $(1^k)\not\subseteq(k-1)\times\ell$.

The kernel of $\pi$ is the image in $\rkl$ of $I^{k-1,\ell}$, which (modulo
$I^{k,\ell}$) is spanned by those $\overline{s_\lambda}$ with $\lambda\subseteq
k\times\ell$ but $\lambda\not\subseteq(k-1)\times\ell$, i.e.\ with exactly $k$ rows
($\ell(\lambda)=k$). This proves the second equality in \eqref{eq:kerpi}.

For the first equality, $e_k=s_{(1^k)}$, and the vertical-strip Pieri rule gives in
$\Lambda$
\[
e_k\cdot s_\mu=s_{\mu+(1^k)}\qquad(\ell(\mu)\le k),
\]
the only vertical strip of size $k$ addable to a partition of length $\le k$ being a
full extra first column. Reducing modulo $I^{k,\ell}$: if $\mu\subseteq k\times\ell$
and $\mu_1\le\ell-1$ then $\mu+(1^k)\subseteq k\times\ell$ and has exactly $k$ rows;
if $\mu_1=\ell$ then $\mu+(1^k)$ has first part $\ell+1>\ell$, so
$\overline{s_{\mu+(1^k)}}=0$. Therefore
\[
e_k\cdot\rkl=\operatorname{span}_{\Q}
\{\,\overline{s_{\mu+(1^k)}}\ :\ \mu\subseteq k\times(\ell-1)\,\}
=\operatorname{span}_{\Q}\{\,\overline{s_\lambda}\ :\ \lambda\subseteq k\times\ell,\ \ell(\lambda)=k\,\},
\]
the last equality by the bijection $\lambda\mapsto\mu=\lambda-(1^k)$ between
$k$-row partitions in $k\times\ell$ and partitions in $k\times(\ell-1)$. This is
exactly $\ker\pi$, proving \eqref{eq:kerpi}. Finally, since $\overline{s_{\mu}}\mapsto
\overline{s_{\mu+(1^k)}}$ raises degree by $k$ and is a bijection onto a basis of
$\ker\pi$, \eqref{eq:kerHilb} follows from $\Hilb(R^{k,\ell-1};q)=\qbin{k-1+\ell}{k}$
(equation \eqref{eq:HilbR} with $\ell$ replaced by $\ell-1$).
\end{proof}

\begin{corollary}[$\pi$ respects the filtration]\label{cor:pifiltr}
For $0\le m\le \ell$ and $k\ge2$ we have $\pi(\Rklm)=R^{k-1,\ell,m}$. Consequently
$\pi$ induces a surjection of graded vector spaces
\[
\overline{\pi}\colon \Rklm/R^{k,\ell,m-1}\twoheadrightarrow
R^{k-1,\ell,m}/R^{k-1,\ell,m-1},
\]
and therefore
\begin{equation}\label{eq:Qdiff}
Q_{k,\ell,m}(q)-Q_{k-1,\ell,m}(q)=\Hilb\big(\ker\overline{\pi};q\big),
\qquad
\ker\overline{\pi}=\frac{\pi^{-1}(R^{k-1,\ell,m-1})\cap \Rklm}{R^{k,\ell,m-1}} .
\end{equation}
\end{corollary}

\begin{proof}
For the nontrivial step we have $m\le k-1$, so $e_1,\dots,e_m$ are among the
generators $e_1,\dots,e_{k-1}$ that $\pi$ sends to the corresponding generators of
$R^{k-1,\ell}$. As $\pi$ is an algebra homomorphism,
$\pi\big(\Q\langle e_1,\dots,e_m\rangle\big)=\Q\langle e_1,\dots,e_m\rangle$ inside
$R^{k-1,\ell}$, i.e.\ $\pi(\Rklm)=R^{k-1,\ell,m}$, and likewise
$\pi(R^{k,\ell,m-1})=R^{k-1,\ell,m-1}$. Thus $\pi$ carries the pair
$(R^{k,\ell,m-1}\subseteq\Rklm)$ onto $(R^{k-1,\ell,m-1}\subseteq R^{k-1,\ell,m})$
and induces a surjection $\overline{\pi}$ on the quotients. For a surjective map of
graded vector spaces, $\Hilb(\text{domain})=\Hilb(\text{image})+\Hilb(\ker)
=\Hilb(\text{codomain})+\Hilb(\ker)$, which is the first identity in \eqref{eq:Qdiff};
the description of $\ker\overline\pi$ is immediate from its definition.
\end{proof}

By Corollary~\ref{cor:pifiltr}, the inductive step \eqref{eq:LHSrec} is equivalent
to the following Hilbert-series identity for the kernel of $\overline\pi$.

\begin{proposition}[Kernel of $\overline\pi$]\label{prop:kernel}
For $1\le m\le\ell$ and $m<k$,
\begin{equation}\label{eq:kernelHilb}
\Hilb\big(\ker\overline\pi;q\big)
= q^{\,m+(k-m)(\ell-m+1)}\,\qbin{\ell}{m}\,\qbin{k-1}{m-1}.
\end{equation}
\end{proposition}

\section{Reduction to two Hilbert-series identities}

The two statements that complete the proof are Proposition~\ref{prop:kernel}
(equivalent to the inductive step \eqref{eq:LHSrec}) and the base case
\eqref{eq:LHSbase}. We record their precise content and the structural reductions
established.

\subsection*{Reformulation of Proposition~\ref{prop:kernel}}
By \eqref{eq:Qdiff} and $\ker\pi=e_k\rkl$ (Lemma~\ref{lem:pi}),
\[
\ker\overline\pi\ \cong\
\frac{\{x\in\Rklm:\ \pi(x)\in R^{k-1,\ell,m-1}\}}{R^{k,\ell,m-1}} .
\]
The target series in \eqref{eq:kernelHilb} factors as
$q^{\,m+(k-m)(\ell-m+1)}\cdot\qbin{\ell}{m}\cdot\qbin{k-1}{m-1}$, where
$\qbin{\ell}{m}=\sum_{\nu\subseteq m\times(\ell-m)}q^{|\nu|}$ is the Hilbert series
of partitions in an $m\times(\ell-m)$ box and
$\qbin{k-1}{m-1}=\sum_{\rho\subseteq(m-1)\times(k-m)}q^{|\rho|}$ that of partitions
in an $(m-1)\times(k-m)$ box, the whole shifted in degree by $m+(k-m)(\ell-m+1)$.

\subsection*{Reformulation of the base case \eqref{eq:LHSbase}}
For $k=m$, $\Rklm=R^{m,\ell,m}=R^{m,\ell}$ (Lemma~\ref{lem:trivial}), so
\[
Q_{m,\ell,m}(q)=\qbin{m+\ell}{m}-\Hilb(R^{m,\ell,m-1};q),
\]
and \eqref{eq:LHSbase} is equivalent to
\begin{equation}\label{eq:basecum}
\Hilb(R^{m,\ell,m-1};q)=\qbin{m+\ell}{m}-q^{m}\qbin{\ell}{m}.
\end{equation}
Here $q^{m}\qbin{\ell}{m}$ is the generating function of partitions with exactly $m$
parts, each in $\{1,\dots,\ell-m+1\}$; equivalently, \eqref{eq:LHSbase} asserts that
$\dim_\Q\big(R^{m,\ell}/R^{m,\ell,m-1}\big)_d$ equals the number of partitions of $d$
with exactly $m$ parts, all $\le\ell-m+1$.

\subsection*{Assembling the induction}
Granting Proposition~\ref{prop:kernel} and the base case \eqref{eq:LHSbase}, the main
identity \eqref{eq:main} follows. Fix $\ell$ and $1\le m\le\ell$ and induct on
$k\ge m$. For $k=m$, \eqref{eq:LHSbase} and \eqref{eq:RHSbase} give
$Q_{m,\ell,m}=q^m\qbin{\ell}{m}=F_{m,\ell,m}$. For $k>m$,
Corollary~\ref{cor:pifiltr} and Proposition~\ref{prop:kernel} give
\[
Q_{k,\ell,m}=Q_{k-1,\ell,m}+q^{\,m+(k-m)(\ell-m+1)}\qbin{\ell}{m}\qbin{k-1}{m-1}
=F_{k-1,\ell,m}+\big(F_{k,\ell,m}-F_{k-1,\ell,m}\big)=F_{k,\ell,m},
\]
using the inductive hypothesis $Q_{k-1,\ell,m}=F_{k-1,\ell,m}$ and
Lemma~\ref{lem:RHSrec}. This completes the induction, hence \eqref{eq:main}, for all
$k\ge m$; for $k<m$ the claim is the degenerate case where $\Rklm=\rkl$ and
\eqref{eq:main} reduces to $m=\ell=k$ being excluded, so the stated range is covered.

\section*{Remaining gap}

The reduction above is rigorous and complete: \eqref{eq:main} is equivalent to the
pair \eqref{eq:LHSrec}--\eqref{eq:LHSbase}; the surjection $\pi$ and its kernel
$e_k\rkl$ are established (Lemma~\ref{lem:pi}); $\pi$ respects the filtration
(Corollary~\ref{cor:pifiltr}); the right-hand side recursion is proved
(Lemma~\ref{lem:RHSrec}); and the cases $m\ge k$, $m=\ell$ are handled
(Lemma~\ref{lem:trivial}). We have additionally shown
(Section~\ref{sec:Fm}) that the natural column-height span $F_m$ is \emph{not} a
viable route: $\Rklm\not\subseteq F_m$ in general (Lemma~\ref{lem:Fmfails}), so the
subalgebra has \emph{no} Schur-class basis cut out by a column-height bound, and the
inductive route through $\pi$ is essential.

What is not yet derived from first principles is the
Hilbert series \eqref{eq:kernelHilb} of $\ker\overline\pi$ (the inductive step) and
the complementary count \eqref{eq:basecum} (the base case $k=m$). Both have been
verified symbolically for all $1\le k,\ell\le5$, $1\le m\le\ell$
(Remark~\ref{rem:verif}); supplying these two computations from the algebra of
$\rkl$ completes the proof of \eqref{eq:main}.

\end{document}
